\textbf{T3: Crossed wires (10 pts)}

a) (1 pt) Current flows through an infinite, straight, thin wire. There is an
externally imposed, uniform magnetic field directed parallel to the wire.
Qualitatively sketch one of the magnetic field lines.

b) (5 pts) Now consider two infinite, straight, thin wires (wires $X$ and $Y$),
each carrying a current $I$ as shown in the figure. The $x$-axis coincides with
wire $X$, while wire $Y$ is parallel to the $y$-axis and passes through the
point $(0,0,-a)$. Let P be the point $(3a,0,r)$. Assuming $r \ll a$, calculate
$d$, the distance of closest approach of the magnetic field line that passes
through P to the wire $X$.

c) (4 pts) Let $L$ be the length of this field line between P and its point of
closest approach to wire $X$. Using values $a = 10\,\mathrm{cm}$ and
$r = 1.0\,\mathrm{mm}$, calculate $L$ to within 20\% relative error.

\begin{center}\includegraphics[width=0.55\linewidth]{figures/EuPhO_2025_Q3_fig1.png}\end{center}
